\section{Target-Local Adaptive C2F Selection}
\label{sec:proposed_algorithm}

\subsection{Prediction-Only Fusion Placement}

The outer stage assigns each target to the UAV nearest its predicted position:
\begin{equation}
	f_q=\arg\min_{m\in\mathcal M}
	\|\widehat{\mathbf p}_q-\mathbf p_m^{\mathrm{UAV}}\|_2 .
	\label{eq:nearest_target_fusion}
\end{equation}
The rule uses neither the current target truth nor its realized RCS. It is a
low-complexity placement heuristic, not a jointly optimal fusion solution.
Given $f_q$, the feasible observation set is
\begin{equation}
\begin{aligned}
	\mathcal E_q=\{(i,j,q):{}&(i,j,q)\text{ is DD-valid},\\
	&j=f_q\text{ or }(j,f_q)\text{ is feasible}\}.
\end{aligned}
	\label{eq:feasible_candidate_set}
\end{equation}

\noindent\textit{Local-delivery bound.}
Under \eqref{eq:received_moments}, the single-observation deflection is
$d(\chi)=\chi^2\delta^2/[v_0(a+(1-a)\chi)]$. Its derivative is nonnegative
for $a\ge0$, so local delivery maximizes this quantity for the same observation.
This does not imply monotonicity of the approximate $P_D$ or optimality of
nearest-target placement when other report links change.

For fixed observation laws and hypothesis-independent replacement, lower
reliability is a random degradation of higher reliability. Optimal detection
at fixed false alarm is therefore nondecreasing with reliability
~\cite{blackwell_experiments}. The supplementary proof concerns the optimal
ROC envelope, not the implemented linear detector.

For a fixed $\mathcal S_q$, let $a_{qm}$ count selected observations received
at UAV $m$. Every fusion choice obeys
$N_{{\rm rem},q}\ge |\mathcal S_q|-\max_m a_{qm}$.
This follows because only observations produced at the fusion UAV avoid a
packet. A feasible nearest-target assignment attaining the bound is
report-minimal for that fixed set; no joint selection optimum is implied.

\subsection{Communication-Constrained Selection}

Let $\widehat{\mathbf P}_D(\mathcal S)$ denote the detector-predicted target
probabilities from \eqref{eq:detector_prediction}. We cap each target probability
at the design requirement and define a smooth weak-target utility
\begin{equation}
\begin{aligned}
	U(\widehat{\mathbf P}_D)=&-Q\tau_\alpha\log\sum_q
	\exp[-\bar P_{D,q}/\tau_\alpha]\\
	&-\frac{\mu_\alpha}{2P_D^{\rm req}}
	\sum_q[P_D^{\rm req}-\widehat P_{D,q}]_+^2,
\end{aligned}
	\label{eq:fair_sensing_utility}
\end{equation}
where $\bar P_{D,q}=\min(\widehat P_{D,q},P_D^{\rm req})$. The first term
prioritizes weak targets, while the second penalizes unresolved detection gaps.
The selector heuristically seeks a high value of the capped objective:
\begin{equation}
\begin{aligned}
	\max_{\mathcal S\subseteq\cup_q\mathcal E_q}\quad
	&F(\mathcal S)=U(\widehat{\mathbf P}_D(\mathcal S))
	-\lambda_c\sum_{e\in\mathcal S}c_e,\\
	\mathrm{s.t.}\quad&|\mathcal S_q|\le L_{\max},\quad
	|\mathcal S|\le L_{\mathrm{tot}},
\end{aligned}
	\label{eq:task_objective}
\end{equation}
where $c_e=10^3\mathbf1\{j\ne f_q\}n/B_{\mathrm c}$ is measured in
milliseconds, matching the implemented price $\lambda_c$. Fixed blocklength makes
$c_e$ a remote-report cardinality price rather than a link-dependent delay.
The caps limit observations, including local ones; the price does not enforce
a prescribed hard communication budget. They imply payload at most
$kL_{\rm tot}$ and serial delay at most $nL_{\rm tot}/B_{\rm c}$.
The greedy replay repeatedly commits the feasible candidate with the largest
positive marginal $F(\mathcal S\cup\{e\})-F(\mathcal S)$. Because the objective
can be non-monotone and the constraints combine target and total caps, we claim
neither global optimality nor a universal submodular approximation ratio.

\subsection{Adaptive C2F Acceleration}
\begin{table}[t]
\caption{Prediction-only adaptive C2F (released V1)}
\label{alg:adaptive_c2f}
\footnotesize
\begin{algorithmic}[1]
\REQUIRE Belief-side tables, $f_q$, caps, $J=25$, $W=1$
\STATE Build triplet sets $\mathcal E_q$; set $S=\varnothing$, $H_q=\varnothing$.
\WHILE{$|S|<L_{\rm tot}$}
\STATE Compute coarse $\Delta F(e\mid S)$ for unselected candidates
satisfying caps and $\Delta D_q>0$.
\STATE For each target, add its highest positive-gain candidate outside
$H_q$ if $|H_q|<\min(J,|\mathcal E_q|)$.
\STATE Let $e^*$ have largest positive gain; stop if none exists.
\STATE Commit $e^*$ to $S$. Insert it into $H_q$ if absent,
replacing the last entry when full.
\ENDWHILE
\STATE Evaluate $\eta^{\rm loc}$ only on $H=\bigcup_q H_q$.
\STATE Reset $S=\varnothing$; greedily add the largest positive fine
$\Delta F$ on $H$, checking caps and $\Delta D_q>0$ each time.
\RETURN $S$ when no admissible positive gain remains or the total cap binds.
\end{algorithmic}
\end{table}


The pseudocode in Table~\ref{alg:adaptive_c2f} constructs a per-target frontier,
then replays the same objective on its fine estimates. The per-target
shortlist cap is $J=25$ and the observation cap is six. A positive-deflection
filter uses $D_q=\sum_{e\in S_q}d(\chi_e)$ under independence.
Cached moments are invalidated when their target's selected set changes.
Ties retain the first candidate in target/transmitter/receiver order.

Let $E$ be the number of feasible candidates and $r$ the number refined. If a
coarse evaluation costs $C_{\mathrm c}$ and a fine evaluation costs
$C_{\mathrm f}$, full refinement costs $EC_{\mathrm f}$, whereas adaptive C2F
has DD evaluation cost approximately
\begin{equation}
	EC_{\mathrm c}+rC_{\mathrm f},\qquad r\le E.
	\label{eq:c2f_complexity}
\end{equation}
This excludes greedy scoring and frontier maintenance. With at most
$L_{\mathrm{tot}}$ commits per stage, scanning requires
$O(L_{\mathrm{tot}}(E+r))$ candidate-score evaluations; a direct implementation costs
$O(L_{\rm tot}(E+r)(L_{\max}+Q))$ for moments and global utility evaluation,
before caching. Nearest placement costs $O(QM)$ and the shortlist stores
at most $QJ$ triplets. When fine DD refinement dominates, the fractional reduction approaches $1-r/E$.
The main experiment used 61.3 refinements for 860.2 feasible candidates
(92.9\% fewer). This is a work-count ratio, not a wall-clock speedup.
No bound on omitted candidates is available, so C2F is not claimed lossless.
